\chapter{Automated Laboratories and Materials Discovery}
\label{ch:07}

If the generation of proposals is cheap and their verification is expensive, then the economic center of gravity of an accelerated science lies in the instruments of verification, and among these the automated laboratory is the one that has most changed in recent years. An automated laboratory in the relevant sense is a facility in which synthesis, processing, and measurement are performed by robotic equipment under computer control, so that a sequence of experiments can be run with little human intervention and results can be fed back to a program that decides what to try next. The conception is not new. Robotic liquid handling and high-throughput screening have long been standard in pharmaceutical research, and the closed-loop selection of experiments by statistical models is established in fields as varied as catalysis and protein engineering. What distinguishes the present prospect is the combination of this infrastructure with generative and predictive models of much greater capacity, so that the loop from hypothesis to measurement to revised hypothesis could in principle operate at a rate and over a space of possibilities that no group of human experimenters could match. The book accepts that this prospect is serious, and it asks what conditions are required for it to produce knowledge and not only data. Recent computational work, such as the structure prediction of Jumper and colleagues~\cite{jumper2021} and the large-scale materials screening of Merchant and colleagues~\cite{merchant2023}, shows how fast candidates can now be generated, and it is cited here as an example of the capability and not as support for any proposal of this book.

The first condition is that the physical operations of the laboratory be reliable and documented. A robotic platform performs the same operation repeatedly, and this is both its strength and its danger. If a pipetting step is miscalibrated, or if a reagent has degraded, or if a furnace profile differs from its nominal setting, every experiment that depends on the faulty step is affected in the same way, and the resulting errors are systematic and not random. Systematic errors are not reduced by repetition and they are not detected by statistical analysis of the replicated data, because the replicates agree with one another. The standard remedies belong to metrology: regular calibration against reference standards, inclusion of control samples with known outcomes in every batch, recording of the provenance of materials and the conditions of each step, and periodic replication of results in a different laboratory with different equipment. A facility that cannot show these records has not produced results that others can use, whatever the sophistication of its software. For this reason the book regards the metrological infrastructure of automated laboratories as part of the verification capacity that must grow in step with generation, and as a proper object of public investment.

The second condition concerns the relation between the model that selects experiments and the experiments themselves. A selection rule is a policy for choosing the next experiment given the results so far, and the policies in common use balance two aims: exploiting regions of the search space in which the results are already good, and exploring regions in which the model is uncertain. The balance is adjustable and its adjustment is a consequential decision, since a policy that exploits too aggressively will converge on a local optimum and report it as the best that can be found, whereas a policy that explores too widely will spend its budget on regions of little promise. The most commonly used criterion, the expected improvement over the best result so far, offers a principled compromise, and a worked comparison, given in the formalization, shows that a candidate with a poorer predicted mean but much greater uncertainty can have a higher expected improvement than a candidate with a better mean and little uncertainty. This is a desirable property of the criterion, but it carries a consequence for how results are reported: the sequence of experiments selected by an adaptive policy is not a random sample of the space, and statements about the distribution of properties over the space, or about the fraction of candidates that succeed, must be corrected for the selection.

The third condition is that claims of discovery be evaluated with attention to multiplicity. A laboratory that tests many thousands of candidates will find, among them, a number that appear to meet the criterion by chance, and if each candidate is judged by an unadjusted significance threshold the list of apparent discoveries will contain many false ones. The statistical literature offers procedures that control the proportion of false discoveries among those reported, of which the procedure of Benjamini and Hochberg~\cite{benjamini1995} is the best known. Control of the false discovery rate is not a substitute for replication, but it provides a defensible basis for deciding which of many candidates deserve the more expensive follow-up. The reporting of a discovery should include the number of candidates tested, the rule by which the reported ones were chosen, and the estimated rate of false discoveries, and the omission of these quantities should be regarded as a defect in the claim, in the same way that the omission of the number of trials is a defect in a claim about the efficacy of a drug.

The discovery of materials has attracted particular attention, and a measured assessment of the claims made on its behalf illustrates the general lessons. Computational screening has been used to propose very large numbers of candidate inorganic compounds, and robotic synthesis has been used to attempt the preparation of some of them. The difficulties that emerge in practice are instructive. A compound that is predicted to be stable may be stable only relative to the competing phases considered, and the prediction can fail if a competing phase has been omitted. A compound that is stable may nonetheless be unsynthesizable by any available route, because the kinetics of formation favor other products. A compound that is synthesized may differ from the predicted structure in ways that matter for its properties, through disorder, impurities, or off-stoichiometry. And a compound that is synthesized in the desired form may be a laboratory curiosity without economic relevance, because the elements it requires are scarce, or because it cannot be produced in quantity at acceptable cost, or because its properties degrade under operating conditions. The path from a promising candidate to a deployed material has been, historically, a matter of decades, and most of the delay is accounted for by the stages that follow the identification of the candidate. An increase in the rate of candidate identification does not change this fact, and the book therefore treats claims of accelerated materials discovery as claims about the earliest stage of a long chain, to be assessed by their effect on the end of the chain.

None of this diminishes the value of the instrument where it is used appropriately. Within a defined domain, such as the optimization of a known family of catalysts or the exploration of the composition space of an established class of alloys, an automated laboratory working under an adaptive policy can find optima more efficiently than manual methods, can detect interactions among variables that would have been overlooked, and can produce datasets of a size and consistency that are themselves a resource for the community. The appropriate governance of such facilities follows from their character as shared infrastructure. They should be open to external users on stated terms, their data should be published with the metadata needed for reuse, their protocols should be recorded in a form that allows other laboratories to reproduce them, and their funding should not depend on the announcement of discoveries, since a dependence of that kind creates an incentive to report optimistic results. These recommendations are the ordinary norms of good science applied to a new kind of facility, and the reason for stating them is that the novelty of the facility and the enthusiasm that surrounds it make it easy to waive them.

\section*{Formalization}

Let a response surface $f:\mathcal X\to\R$ be unknown, and let a model provide for each candidate $x$ a Gaussian predictive distribution $f(x)\sim\mathcal N(\mu(x),\sigma^2(x))$. Let $f^\ast$ be the best value observed so far, with larger values preferred. The \emph{expected improvement} of $x$ is $\mathrm{EI}(x)=\E[\max(f(x)-f^\ast,0)]$.

\begin{proposition}
With $z=(\mu-f^\ast)/\sigma$, and writing $\Phi$ and $\phi$ for the standard normal distribution and density functions, $\mathrm{EI}(x)=\sigma\bigl[z\Phi(z)+\phi(z)\bigr]$. The function $\mathrm{EI}$ is strictly increasing in $\mu$ for fixed $\sigma>0$ and strictly increasing in $\sigma$ for fixed $\mu$.
\end{proposition}

\begin{proof}
Write $f=\mu+\sigma\xi$ with $\xi$ standard normal. Then $\E[(\mu-f^\ast+\sigma\xi)^+]=\sigma\,\E[(z+\xi)^+]=\sigma\int_{-z}^{\infty}(z+u)\phi(u)\,du=\sigma[z\Phi(z)+\phi(z)]$, using $\int_{-z}^\infty u\phi(u)\,du=\phi(z)$. Differentiating, $\partial\mathrm{EI}/\partial\mu=\Phi(z)>0$ and $\partial\mathrm{EI}/\partial\sigma=\phi(z)>0$.
\end{proof}

The comparison cited in the text is as follows. With incumbent $f^\ast=1$, a candidate A with $\mu=0.5$ and $\sigma=1.0$ has $z=-0.5$ and $\mathrm{EI}\approx0.198$, while a candidate B with $\mu=0.9$ and $\sigma=0.1$ has $z=-1$ and $\mathrm{EI}\approx0.0083$. The candidate whose predicted mean is lower by $0.4$ is selected, by a factor of about twenty-four in expected improvement, because its uncertainty is ten times greater. This is exploration, and the same calculation shows why the selected experiments are not a random sample.

For multiplicity, let $m$ hypotheses be tested with $p$-values $p_1,\dots,p_m$, of which $m_0$ are true nulls. Order them $p_{(1)}\le\dots\le p_{(m)}$ and, for a target level $\alpha\in(0,1)$, let $\hat k=\max\{i:p_{(i)}\le i\alpha/m\}$, rejecting the hypotheses with the $\hat k$ smallest $p$-values (and none if no such $i$ exists). The false discovery rate is $\mathrm{FDR}=\E[V/\max(R,1)]$, where $R$ is the number of rejections and $V$ the number of true nulls among them.

\begin{proposition}[Benjamini and Hochberg, 1995]
If the $p$-values of the true nulls are independent of one another and of the non-null $p$-values, and each is uniformly distributed on $[0,1]$, then $\mathrm{FDR}\le\alpha\,m_0/m\le\alpha$.
\end{proposition}

The result is cited and not reproved. In the screening of $m=1000$ candidates with $m_0=900$ true nulls at $\alpha=0.1$ it guarantees $\mathrm{FDR}\le0.09$: of the candidates passed to follow-up, no more than nine per cent are expected to be false leads. The bound holds only under the stated independence, which is itself a modelling assumption about the generating process, and the bound should be reported with the check that supports the assumption, or replaced by a procedure valid under dependence.
